\setcounter{section}{56} % tema número N-1
\section{Multi Usuario}

Nombre completo: Tipos de sistemas de información multiusuario. Sistemas grandes, medios y pequeños. Servidores de datos y aplicaciones. Virtualización de servidores.

\localtableofcontents

\subsection{Sistemas de información multiusuario}
La clasificación es aproximada, y tiene herencias legacy. De mayor a menor son

\begin{description}
    \item[Mainframes] Soportan una elevada carga de trabajo y usuarios. Son estructuras centralizadas. Tienen una alta fiabilidad y se diseñan para trabajar sin interrupción durante años
    \item[Super ordenadores] Gran potencia en base a arquitecturas paralelas. Se gestionan por personal especializado.
    \item[Miniordenador o Servidor] Ordenador medio. Puede realizar cálculos complejos y soporta una gran cantidad de usuarios. Están a mitad de camino entre el mainframe y el microordenador.
    \item[Estaciones de trabajo] Ordenadores personales con altas prestaciones. Se usan comúnmente para CAD o diseño 3D
    \item[Microordenadores] Ordenadores de sobremesa o portátiles.
    \item[Otros] Tablets y smartphones.
\end{description}

\subsection{Servidores de datos y aplicaciones}
Según el papel que asumen se pueden clasificar como \textbf{dedicados} si se usan para un único fin. o \textbf{no dedicados} si aceptan varios servicios simultáneamente.

Los usos mas comunes son
\begin{description}
    \item[web] se encargan de la capa de presentación
    \item[aplicaciones] Ejecutan la lógica de negocio
    \item[datos] Sistemas de base de datos
    \item[correo]
    \item[multimedia] Distribuye contenido multimedia
    \item[noticias]
    \item[colaboración] Provee de herramientas colavorativas para sincronizar el trabajo de varios usuarios
    \item[proxy] Intermediario de conexión para sus usuarios
\end{description}

\subsection{Virtualización}
Gracias a tecnologías como la nube se permite una reducción de costes operativos. Cambiando el paradigma a los servicios por uso en lugar de por disponibiilidad. Gracias a ello se pueden hacer despliegues automáticos, escalado bajo demanda o continuidad de negocio.

Según cuanto se virtualiza y qué capa queda expuesta, la virtualización puede ser de varios tipos.
\begin{description}
    \item[Emulación] se emulta completamente un tipo de arquitectura en otra. (DosBox)
    \item[Nativa o total] La máquina virtual emula una cantidad suficiente de hardware para que el entorno virtualizado funcione completamente sin modificar (VMWare, VirtualBox)
    \item[para-virtualización] El host no emula el hardware, sino que ofrece una serie de APIs a un sistema convenientemente modificado (UML, xVM)
    \item[Virtualización a nivel de SO] Los huéspedes comparten el mismo kernel que el host. Creandose instancias independientes entre si. (Linux V Server, OpenVZ)
    \item[Virtualización de aplicaciones] Las aplicaciones tienen su propio entorno virtualizado con lo necesario para ejecutarse tanto sobre un servidor, como un cliente.
\end{description}